\chapter{Raccordement torsionnel d'une région piégée et construction d'une branche cosmologique descendante}
\label{chap:universo-descendiente}

\section{Raccordement torsionnel d'une région piégée : strates exacte, conditionnée et physique de la branche cosmologique descendante}

Le \cref{ch:rebote-ec} établit la loi de mémoire
\begin{equation}
 M_n=\left(\frac{a_n}{a_0}\right)^{-6}
 \label{eq:descendant-memory-law}
\end{equation}
et formule sa réalisation par une densité quadratique de spin dans
Einstein--Cartan. Le présent chapitre étudie le passage supplémentaire d'un
rebond torsionnel à une région causale expansive située au-delà de l'intérieur
piégé d'un trou noir.

L'exposé distingue \(3\) strates :
\begin{enumerate}
 \item \emph{structure exacte} : orientation TRIT, involution de Möbius,
       conservation isométrique du canal et loi \(V^{-2}\) ;
 \item \emph{conséquence conditionnée} : rebond d'Einstein--Cartan,
       raccordement géométrique et ouverture causale sous hypothèses déclarées ;
 \item \emph{réalisation géométrique classifiée} : existence explicite de
       la continuation descendante minimale et obstruction au cobordisme
       ramifié classique avec changement topologique.
\end{enumerate}
Cette stratification permet de développer l'image complète et de conserver le
statut de chaque affirmation.

L'antécédent commun est un état enrichi
\begin{equation}
 X^{\rm enr}_{\TPK}
 =(x,\text{fibre},\text{orientation},\text{carry},
   \text{frontière},\text{route},\text{mémoire},\text{observateur}),
 \label{eq:descendant-enriched-state}
\end{equation}
produit par \(\mathrm{APP}\to\TRIT\to\TPK\). Les lecteurs causal,
torsionnel, topologique et quantique s'appliquent à ce même état et
conservent sa généalogie.

\section{Orientation causale des \(3\) feuillets}

La grammaire de TRIT se réalise au moyen des algèbres
\begin{equation}
 \mathbb A_\tau=\R[X]/(X^2+\tau),
 \qquad
 J_\tau^2=-\tau I,
 \qquad
 \tau\in\{+1,0,-1\}.
 \label{eq:descendant-TRIT-algebras}
\end{equation}
Leurs exponentielles sont
\begin{align}
 \ee^{tJ_{+1}}&=\cos t\,I+\sin t\,J_{+1},
 \label{eq:descendant-elliptic}\\
 \ee^{tJ_0}&=I+tJ_0,
 \label{eq:descendant-parabolic}\\
 \ee^{tJ_{-1}}&=\cosh t\,I+\sinh t\,J_{-1}.
 \label{eq:descendant-hyperbolic}
\end{align}
Le feuillet \(+1\) porte la clôture, le retour et la mémoire ; le feuillet \(0\) porte le seuil,
le présent et la première variation ; le feuillet \(-1\) porte l'ouverture, la propagation et le
futur two-way.

Soit \(\vartheta\) un scalaire d'expansion orienté. Nous définissons le lecteur
causal
\begin{equation}
 \boxed{
 \tau_{\rm cau}(\vartheta):=-\operatorname{sgn}(\vartheta)
 \in\{-1,0,+1\}.}
 \label{eq:descendant-causal-reader}
\end{equation}
Une expansion positive publie l'ouverture \(-1\) ; une expansion nulle
publie le seuil \(0\) ; une expansion négative publie la clôture \(+1\).

Dans une famille de sphères, \(\vartheta=\theta_{(\ell)}\) peut être pris comme
l'expansion nulle sortante. Dans une région homogène ou anisotrope,
\(\vartheta=\Theta\) peut être pris comme l'expansion volumique. Ces
deux observables occupent des domaines distincts et partagent la même règle de
signe.

\begin{proposition}[Séquence causale à double seuil]
\label{prop:descendant-double-threshold}
Supposons les conditions de signe
\begin{equation}
 \theta_{(\ell)}>0\ \text{à l'extérieur},
 \quad
 \theta_{(\ell)}=0\ \text{à l'horizon},
 \quad
 \theta_{(\ell)}<0\ \text{dans l'intérieur piégé},
 \label{eq:descendant-null-signs}
\end{equation}
et
\begin{equation}
 \Theta<0\ \text{avant le rebond},
 \quad
 \Theta=0\ \text{au rebond},
 \quad
 \Theta>0\ \text{après le rebond}.
 \label{eq:descendant-volume-signs}
\end{equation}
Le lecteur \cref{eq:descendant-causal-reader} produit la séquence
\begin{equation}
 \boxed{
 -1_{\rm exterior}
 \longrightarrow0_{\rm horizonte}
 \longrightarrow+1_{\rm interior}
 \longrightarrow0_{\rm rebote}
 \longrightarrow-1_{\rm descendiente}.}
 \label{eq:descendant-five-stage-sequence}
\end{equation}
\end{proposition}

\begin{proof}
L'affirmation résulte de l'application de
\(\tau_{\rm cau}(\vartheta)=-\operatorname{sgn}(\vartheta)\) à chaque signe de
\cref{eq:descendant-null-signs,eq:descendant-volume-signs}.
\end{proof}

Les deux états \(0\) sont des réalisations du même archétype de seuil. Le
premier exprime la marginalité nulle ; le second exprime la marginalité
volumique. L'identité du signe conserve la grammaire, tandis que l'
observable déclaré détermine le contenu physique de chaque seuil.

\section{Intérieur anisotrope et mémoire quadratique}

Une région intérieure à symétrie sphérique admet la réduction de
Kantowski--Sachs
\begin{equation}
 \dd s^2=-N(\tau)^2\dd\tau^2
 +a_\parallel(\tau)^2\dd\chi^2
 +a_\perp(\tau)^2\dd\Omega_2^2.
 \label{eq:descendant-KS-metric}
\end{equation}
Le volume d'une cellule comobile est
\begin{equation}
 V(\tau)=V_0a_\parallel a_\perp^2,
 \label{eq:descendant-KS-volume}
\end{equation}
et son expansion est
\begin{equation}
 \Theta
 :=\frac1N\frac{\dot V}{V}
 =H_\parallel+2H_\perp,
 \qquad
 H_i:=\frac1N\frac{\dot a_i}{a_i}.
 \label{eq:descendant-KS-expansion}
\end{equation}

Une charge comobile de spin \(N_s\) produit une densité
\begin{equation}
 n_s=\frac{N_s}{V},
 \label{eq:descendant-spin-number-density}
\end{equation}
et l'équation de Cartan, linéaire en le courant, engendre une correction
quadratique
\begin{equation}
 \rho_s=s_0\left(\frac V{V_0}\right)^{-2}.
 \label{eq:descendant-spin-density}
\end{equation}
La branche de Perron publie exactement la même covariance :
\begin{equation}
 M_n=V_n^{-2}.
 \label{eq:descendant-Perron-volume}
\end{equation}
Le transducteur torsionnel
\begin{equation}
 \mathcal T_{\rm EC}:M_n\longmapsto
 \rho_{s,n}=s_0M_n
 \label{eq:descendant-EC-transducer}
\end{equation}
conserve la positivité et l'exposant de volume. L'amplitude \(s_0\) provient
du courant de Cartan et maintient son obligation de sélection
prospective.

Pour la discussion locale, nous adoptons les unités géométriques \(c=1\). Une
contrainte effective de Kantowski--Sachs a la forme
\begin{equation}
 \frac{\Theta^2}{3}
 =8\pi G(\rho_m-\rho_s)
 +\sigma^2-\frac12\,{}^{(3)}R+\Lambda,
 \label{eq:descendant-KS-constraint}
\end{equation}
où \(\sigma^2\) est le cisaillement et \({}^{(3)}R\) la courbure du feuillet.
Un rebond volumique régulier satisfait
\begin{equation}
 \boxed{
 \Theta(\tau_b)=0,
 \qquad
 \frac1N\dot\Theta(\tau_b)>0,}
 \label{eq:descendant-bounce-conditions}
\end{equation}
avec la finitude positive de \(a_\parallel,a_\perp\) et des invariants
de courbure et de torsion.

\begin{criterion}[Dominance torsionnelle anisotrope]
\label{crit:descendant-torsion-dominance}
Supposons que, pendant la contraction,
\begin{equation}
 \rho_s=s_0V^{-2},
 \qquad
 \sigma^2=\Sigma_0^2V^{-2}+o(V^{-2}),
 \label{eq:descendant-torsion-shear-scaling}
\end{equation}
et que la matière, la courbure et \(\Lambda\) croissent avec une puissance strictement
inférieure à \(V^{-2}\). L'inégalité
\begin{equation}
 \boxed{8\pi Gs_0>\Sigma_0^2}
 \label{eq:descendant-torsion-dominates-shear}
\end{equation}
est une condition suffisante pour que le terme torsionnel domine le
cisaillement dans le régime de petit volume. L'existence d'une racine
transverse de \cref{eq:descendant-KS-constraint} et l'inégalité
\cref{eq:descendant-bounce-conditions} complètent le critère de rebond.
\end{criterion}

La dominance asymptotique apporte le mécanisme ; la racine transverse apporte
l'événement géométrique. Cette distinction conserve la dépendance par rapport à la
courbure, à l'anisotropie, à l'équation d'état et aux données de frontière.

\section{Condition nonadique pour le niveau de rebond}

Écrivons \(x:=a/a_0\) pour l'échelle isotrope adimensionnelle. Dans les
retours complets \(n=9q\), la branche de Perron satisfait
\begin{equation}
 x_{9q}=\frac{a_{9q}}{a_0}=\varphi_{\HMT}^{3q}.
 \label{eq:descendant-nonadic-scale}
\end{equation}
Dans la réduction isotrope avec
\begin{equation}
 \rho_m=\rho_0x^{-3(1+w)},
 \qquad
 \rho_s=s_0x^{-6},
 \qquad
 w<1,
 \label{eq:descendant-isotropic-densities}
\end{equation}
l'égalité des densités fixe
\begin{equation}
 x_b^{3(1-w)}=\frac{s_0}{\rho_0}.
 \label{eq:descendant-isotropic-bounce}
\end{equation}
Si le rebond est exigé sur un tour complet du TPK, l'élimination de
\(x_b\) produit
\begin{equation}
 \boxed{
 q_b=
 \frac{\log(s_0/\rho_0)}
 {9(1-w)\log\varphi_{\HMT}},
 \qquad
 q_b\in\Z.}
 \label{eq:descendant-quantized-bounce}
\end{equation}
L'intégralité de \(q_b\) est une condition supplémentaire de synchronisation,
dérivée de l'hypothèse de coïncidence entre le rebond et le retour
nonadique. Un rebond continu entre deux tours satisfait
\cref{eq:descendant-isotropic-bounce} et possède un autre statut.

\begin{proposition}[Réfutabilité du rebond synchronisé]
\label{prop:descendant-quantization-falsifier}
\((s_0,\rho_0,w)\) étant fixés prospectivement, la réalisation synchronisée
n'existe que si le nombre de
\cref{eq:descendant-quantized-bounce} est un entier. La distance
\begin{equation}
 \delta_9:=\operatorname{dist}(q_b,\Z)
 \label{eq:descendant-quantization-defect}
\end{equation}
est un critère quantitatif de réfutation et s'annule exactement à un niveau permis.
\end{proposition}

\begin{proof}
La condition \(n=9q\) exige \(q\in\Z\). Substituer
\cref{eq:descendant-nonadic-scale} dans
\cref{eq:descendant-isotropic-bounce} et prendre les logarithmes produit la formule.
La caractérisation de \(\delta_9\) découle de la définition de la distance aux
entiers.
\end{proof}

\section{Données de raccordement sur l'hypersurface de rebond}

Soit \(\Sigma_b\) une hypersurface spatiale séparant une région
contractante \(\mathcal M_-\) d'une région expansive
\(\mathcal M_+\). Notons \(h_{ij}\) la métrique induite,
\(n^\mu\) la normale orientée et
\(\mathring K_{ij}^\pm\) les courbures extrinsèques calculées avec la
connexion de Levi--Civita de chaque côté. Nous adoptons
\begin{equation}
 [X]:=X^+\big|_{\Sigma_b}-X^-\big|_{\Sigma_b},
 \qquad
 \kappa:=\frac{8\pi G}{c^4}.
 \label{eq:descendant-jump-kappa}
\end{equation}

Les conditions métriques d'Israel sont
\begin{align}
 [h_{ij}]&=0,
 \label{eq:descendant-Israel-first}\\
 [\mathring K_{ij}-h_{ij}\mathring K]
 &=-\kappa S^{\Sigma}_{ij},
 \label{eq:descendant-Israel-second}
\end{align}
où \(S^\Sigma_{ij}\) est le tenseur énergie--impulsion surfacique dans la
convention de normale fixée. Un raccordement sans couche matérielle satisfait
\begin{equation}
 [h_{ij}]=0,
 \qquad
 [\mathring K_{ij}]=0.
 \label{eq:descendant-shell-free-Israel}
\end{equation}
L'inversion de \(\Theta\) passe alors par un zéro continu et est
réalisée par la dynamique, au lieu d'être introduite par un saut de la
normale.

Dans la formulation du premier ordre, soit
\(\omega^{ab}=\mathring\omega^{ab}+C^{ab}\) la connexion avec contorsion,
et écrivons \(T_\mu:=T^\lambda{}_{\mu\lambda}\) pour la trace de la
torsion.
Nous fixons la normalisation du courant de spin au moyen de l'équation
de Cartan
\begin{equation}
 \mathscr C(T)^{\lambda}{}_{\mu\nu}
 :=T^{\lambda}{}_{\mu\nu}
 +\delta^\lambda_\mu T_\nu
 -\delta^\lambda_\nu T_\mu
 =\kappa s^{\lambda}{}_{\mu\nu}.
 \label{eq:descendant-Cartan-equation}
\end{equation}
Sa partie distributionnelle sur \(\Sigma_b\) exige
\begin{equation}
 \boxed{
 \mathscr C(T_\Sigma)
 =\kappa s_\Sigma.}
 \label{eq:descendant-Cartan-junction}
\end{equation}
En l'absence d'une couche superficielle de spin,
\(s_\Sigma=0\), la connexion admissible est dépourvue de composante torsionnelle
delta et la contorsion tangentielle possède un raccordement régulier. Une action avec
terme de Holst ou couplages fermioniques non minimaux remplace
\(\mathscr C\) par l'opérateur algébrique correspondant ; les conditions
s'obtiennent en intégrant cette équation à travers la couche.

\begin{criterion}[Raccordement Israel--Cartan régulier]
\label{crit:descendant-Israel-Cartan}
Une couture torsionnelle régulière à travers \(\Sigma_b\) requiert :
\begin{enumerate}[label=\roman*)]
 \item la continuité de la costructure induite, équivalente à
       \cref{eq:descendant-Israel-first} à jauge de Lorentz près ;
 \item l'équation d'Israel \cref{eq:descendant-Israel-second} avec toute couche
       superficielle déclarée ;
 \item l'équation de Cartan \cref{eq:descendant-Cartan-junction} avec tout
       courant superficiel déclaré ;
 \item la continuité des charges de jauge et de l'orientation de route ;
 \item la finitude des invariants de courbure et de torsion.
\end{enumerate}
\end{criterion}

La connexion d'Ashtekar--Barbero
\begin{equation}
 A_i^a=\Gamma_i^a+\gamma_{\rm BI}K_i^a
 \label{eq:descendant-Barbero-connection}
\end{equation}
peut agir comme lectrice de la couture lorsque les deux côtés partagent
la costructure, \(\gamma_{\rm BI}\) et la normalisation. Sa continuité est une
conséquence conditionnée des données de raccordement correspondantes ; sa
vérification exige de transporter explicitement \(\Gamma\) et \(K\) vers la
même section de bord.

\begin{proposition}[Involution conjointe de Barbero et de la courbure extrinsèque]
\label{prop:descendant-Barbero-involution}
Considérons l'involution algébrique
\begin{equation}
 \mathfrak J_b:
 (\Gamma_i^a,K_i^a,\gamma_{\rm BI})
 \longmapsto
 (\Gamma_i^a,-K_i^a,-\gamma_{\rm BI}).
 \label{eq:descendant-Barbero-involution}
\end{equation}
Alors
\begin{equation}
 \mathfrak J_b(A_i^a)=A_i^a,
 \label{eq:descendant-Barbero-A-invariant}
\end{equation}
et tout spectre d'aire dépendant de \(|\gamma_{\rm BI}|\) reste
invariant. L'identification de \(\mathfrak J_b\) à l'inversion
dynamique du rebond requiert que l'entrelaceur Israel--Cartan transporte
l'orientation de \(K\) et le feuillet de \(\gamma_{\rm BI}\) conformément à
\cref{eq:descendant-Barbero-involution}.
\end{proposition}

\begin{proof}
Le produit des deux facteurs orientés satisfait
\((-\gamma_{\rm BI})(-K_i^a)=\gamma_{\rm BI}K_i^a\), tandis que \(\Gamma_i^a\)
reste fixe. La dépendance de l'aire par rapport à \(|\gamma_{\rm BI}|\) est paire.
\end{proof}

\section{Horizon nul comme hypersurface marginale et condition initiale de raccordement}

L'hypersurface de l'horizon \(\Sigma_H\) est de nature nulle et requiert des
données de raccordement adaptées. Soient \(q_{AB}\) la métrique bidimensionnelle des
sections, \(\ell^\mu\) le générateur nul et \(N^\mu\) un transversal avec
\(\ell\cdot N=-1\). La courbure transversale est
\begin{equation}
 \mathcal C_{AB}
 :=q_A{}^\mu q_B{}^\nu\nabla_\mu N_\nu.
 \label{eq:descendant-null-transverse-curvature}
\end{equation}
Un raccordement nul sans couche impulsive conserve \(q_{AB}\), la classe de
\(\mathcal C_{AB}\) qui détermine la tension superficielle et le flux de
Cartan tangentiel. La marginalité
\begin{equation}
 \theta_{(\ell)}\big|_{\Sigma_H}=0
 \label{eq:descendant-null-marginality}
\end{equation}
réalise le premier état \(0\) de
\cref{eq:descendant-five-stage-sequence}. La formulation nulle maintient
cette condition séparée de la condition volumique
\(\Theta|_{\Sigma_b}=0\).

\section{Ruban de Möbius et mémoire d'orientation}

Sur les mots dodécaphasiques agit l'involution généalogique
\begin{equation}
 C_M(\mathbf t):=-\operatorname{rev}(\mathbf t),
 \qquad
 C_M^2=I.
 \label{eq:descendant-CM}
\end{equation}
La droite orientée associée s'obtient par le quotient
\begin{equation}
 (\mathbf t,\chi)\sim(C_M\mathbf t,-\chi).
 \label{eq:descendant-Mobius-quotient}
\end{equation}
Lorsque la base est un lacet de retour, cette identification définit un ruban
d'holonomie \(-I\) : un tour inverse l'orientation et deux tours la
rétablissent.

Soit \(\mathscr H_\gamma\) l'holonomie sur un tour. La hiérarchie
\begin{equation}
 \mathcal R_{\rm vec}
 \subseteq\mathcal R_{\rm ray}
 \subseteq\mathcal R_{\rm obs}
 \label{eq:descendant-return-hierarchy}
\end{equation}
distingue retour de vecteur, retour de rayon et retour d'observable. Pour
\(\mathscr H_\gamma=-I\), le rayon revient en un tour, le vecteur orienté
revient en deux et la mémoire de signe conserve le parcours.

\begin{proposition}[Statut topologique du ruban]
\label{prop:descendant-Mobius-status}
La relation \cref{eq:descendant-Mobius-quotient} définit une droite réelle
non orientable sur le lacet exactement lorsque sa première classe de
Stiefel--Whitney satisfait
\begin{equation}
 w_1=1\in H^1(S^1;\Z_2).
 \label{eq:descendant-w1}
\end{equation}
Cette affirmation concerne le fibré généalogique d'orientation. Une
réalisation comme topologie de l'espace-temps requiert un morphisme vers
le fibré des repères ou des spineurs qui conserve la causalité et la structure de
Cartan.
\end{proposition}

\begin{proof}
Une droite réelle sur \(S^1\) est déterminée par la parité de sa
fonction de transition. La transition \(-1\) de
\cref{eq:descendant-Mobius-quotient} représente la classe non triviale ; deux
tours élèvent la transition au carré et produisent \(+1\).
\end{proof}

Le ruban fournit une représentation exacte de l'orientation et de la mémoire.
La continuité physique à travers l'horizon et le rebond se formule
au moyen des raccordements des sections précédentes ; les deux structures sont
coordonnées par un entrelaceur et conservent leurs domaines.

\section{Observateur relatif, transport parallèle et auto-inertie dans la région cosmologique descendante}

L'état dynamique étendu comprend l'observateur et la mémoire :
\begin{equation}
 \Gamma=(x,\mathcal O,m),
 \label{eq:descendant-observer-state}
\end{equation}
et le cadre accessible à l'observateur est une restriction de la connexion totale,
\begin{equation}
 \mathfrak F_{\mathcal O}(\Gamma)
 :=\operatorname{Res}_{\mathcal O}
 (\Gamma,\nabla_{\TPK}).
 \label{eq:descendant-observer-frame}
\end{equation}
Une trajectoire horizontale du système complet,
\begin{equation}
 \nabla^{\rm tot}_{\dot\Gamma}\dot\Gamma=0,
 \label{eq:descendant-total-horizontal}
\end{equation}
peut publier une accélération dans un sous-système \(S\) :
\begin{equation}
 a_S=\nabla^S_{\dot\gamma_S}\dot\gamma_S
 =\mathcal K_{\pi_S}(\dot\Gamma,\dot\Gamma),
 \label{eq:descendant-projected-acceleration}
\end{equation}
où \(\mathcal K_{\pi_S}\) est le second tenseur fondamental de la
projection. L'auto-inertie désigne l'horizontalité de la totalité et la
réponse induite dans ses projections.

Le couplage entre système, appareil et registre se décrit au moyen d'une
isométrie
\begin{equation}
 V_{\rm obs}:\mathcal H_{\rm in}\longrightarrow
 \mathcal H_{\rm sys}\otimes\mathcal H_{\rm mem},
 \qquad
 V_{\rm obs}^*V_{\rm obs}=I.
 \label{eq:descendant-observer-isometry}
\end{equation}
Soit \(\{P_y\}\) la PVM du registre dans \(\mathcal H_{\rm mem}\), avec
\begin{equation}
 P_yP_z=\delta_{yz}P_y,
 \qquad
 \sum_yP_y=I_{\rm mem}.
 \label{eq:descendant-observer-PVM}
\end{equation}
L'isométrie et cette PVM induisent l'instrument complètement positif
\begin{equation}
 \mathcal I_y(\rho)
 :=\Tr_{\rm mem}\!\left[
 (I_{\rm sys}\otimes P_y)V_{\rm obs}\rho V_{\rm obs}^*
 (I_{\rm sys}\otimes P_y)
 \right],
 \qquad
 \sum_y\Tr\mathcal I_y(\rho)=\Tr\rho.
 \label{eq:descendant-observer-instrument}
\end{equation}
Les états relatifs s'obtiennent alors par
\begin{equation}
 p(y)=\Tr\mathcal I_y(\rho),
 \qquad
 \rho_y=\frac{\mathcal I_y(\rho)}{p(y)}.
 \label{eq:descendant-relative-state}
\end{equation}
La section \(\rho_y\) est l'état relatif au registre \(y\) ; l'état
global conserve tous les canaux dans l'image de l'isométrie. Cette
structure formalise la parenté avec Everett au niveau précis des
états relatifs et laisse l'ontologie des branches comme choix physique
ultérieur.

\section{Canal global de la région descendante}

La version opératorielle d'une ouverture descendante utilise une isométrie
\begin{equation}
 V_{\rm pc}:\mathcal H_{\rm in}\longrightarrow
 \mathcal H_{\rm parent}\otimes\mathcal H_{\rm child},
 \qquad
 V_{\rm pc}^*V_{\rm pc}=I.
 \label{eq:descendant-parent-child-isometry}
\end{equation}
Le canal conjoint et ses canaux complémentaires sont
\begin{align}
 \mathcal N_{\rm pc}(\rho)
 &=V_{\rm pc}\rho V_{\rm pc}^*,
 \label{eq:descendant-global-channel}\\
 \mathcal N_{\rm parent}(\rho)
 &=\Tr_{\rm child}\mathcal N_{\rm pc}(\rho),
 \qquad
 \mathcal N_{\rm child}(\rho)
 =\Tr_{\rm parent}\mathcal N_{\rm pc}(\rho).
 \label{eq:descendant-complementary-channels}
\end{align}

\begin{theorem}[Conservation globale et distribution relative]
\label{thm:descendant-global-conservation}
L'application \(\mathcal N_{\rm pc}\) conserve la trace et la pureté d'un état
pur sur son image. Les canaux parentaux et descendants peuvent produire des
états mixtes, et pour une sortie bipartite pure ils satisfont
\begin{equation}
 S(\rho_{\rm parent})=S(\rho_{\rm child}).
 \label{eq:descendant-equal-entanglement}
\end{equation}
L'information globale est contenue dans les sorties et leurs corrélations.
\end{theorem}

\begin{proof}
L'identité \(V_{\rm pc}^*V_{\rm pc}=I\) implique
\[
 \Tr(V_{\rm pc}\rho V_{\rm pc}^*)
 =\Tr(\rho V_{\rm pc}^*V_{\rm pc})=\Tr\rho.
\]
Une isométrie envoie des vecteurs unitaires sur des vecteurs unitaires. L'égalité des
entropies réduites découle de la décomposition de Schmidt de l'état
bipartite pur.
\end{proof}

La conservation d'une charge \(Q\) ou de l'énergie exige en outre les
entrelacements
\begin{equation}
 Q_{\rm out}V_{\rm pc}=V_{\rm pc}Q_{\rm in},
 \qquad
 H_{\rm out}V_{\rm pc}=V_{\rm pc}H_{\rm in}.
 \label{eq:descendant-charge-energy-intertwining}
\end{equation}
L'isométrie conserve la probabilité ; les équations
\cref{eq:descendant-charge-energy-intertwining} conservent les grandeurs
dynamiques déclarées. Une sortie admissible distribue l'information dans des
corrélations et respecte l'impossibilité de cloner un état arbitraire en
deux copies identiques.

La provenance des secteurs \(90/120\) est conservée par l'
hamiltonien holographique modulaire
\begin{equation}
 K_{\rm HHM}=c_0I+\Delta_{\rm mod}
 (90N_{90}+120N_{120}).
 \label{eq:descendant-HHM}
\end{equation}
L'occupation totale et la mémoire pondérée,
\begin{equation}
 N=N_{90}+N_{120},
 \qquad
 D=90N_{90}+120N_{120},
 \label{eq:descendant-HHM-data}
\end{equation}
Le transport de l'occupation totale est imposé comme identité opératorielle
indépendante :
\begin{equation}
 N^{\rm out}V_{\rm pc}=V_{\rm pc}N^{\rm in}.
 \label{eq:descendant-number-intertwining}
\end{equation}
\(N\) et la mémoire pondérée \(D\) étant conservés, les deux secteurs se
reconstruisent de manière univoque :
\begin{equation}
 N_{90}=4N-\frac D{30},
 \qquad
 N_{120}=\frac D{30}-3N.
 \label{eq:descendant-HHM-recovery}
\end{equation}
La condition de transport
\begin{equation}
 K_{\rm HHM}^{\rm out}V_{\rm pc}
 =V_{\rm pc}K_{\rm HHM}^{\rm in}
 \label{eq:descendant-HHM-intertwining}
\end{equation}
Les identités
\cref{eq:descendant-number-intertwining,eq:descendant-HHM-intertwining}
préservent conjointement l'occupation et la provenance sectorielle à travers la
couture. Ces égalités s'ajoutent aux obligations dynamiques d'une
réalisation complète.

\section{Continuation descendante et cobordisme ramifié}

La dénomination \emph{univers descendant} admet deux niveaux
géométriques.

\begin{definition}[Continuation descendante minimale]
\label{def:descendant-minimal}
Une continuation descendante est une région
\(\mathcal M_{\rm child}\) située dans le futur du rebond, avec
\begin{equation}
 \Theta>0,
 \qquad
 \text{courbure et torsion finies},
 \qquad
 \text{une fonction temporelle globale dans son domaine causal}.
 \label{eq:descendant-minimal-conditions}
\end{equation}
La région peut appartenir à un prolongement connexe de l'intérieur et
conserve la topologie des feuillets spatiaux.
\end{definition}

\begin{definition}[Cobordisme parental--descendant]
\label{def:descendant-cobordism}
La version ramifiée est un cobordisme effectif
\begin{equation}
 \mathcal W:\Sigma_{\rm in}\longrightarrow
 \Sigma_{\rm parent}\sqcup\Sigma_{\rm child},
 \label{eq:descendant-cobordism}
\end{equation}
muni d'une structure de Cartan, de données de raccordement
\cref{crit:descendant-Israel-Cartan} et d'un canal
\cref{eq:descendant-parent-child-isometry}.
\end{definition}

La continuation minimale requiert un rebond et une extension causale. Le
cobordisme ramifié ajoute un changement topologique et une séparation des sorties.
Dans une géométrie lorentzienne classique, non dégénérée et globalement
hyperbolique, le changement topologique est soumis à des obstructions causales.
La version ramifiée doit donc se réaliser au moyen d'une région
quantique ou torsionnelle où l'on déclare quelle hypothèse classique est modifiée,
ou par une interprétation de deux domaines causaux au sein d'une même
topologie globale.

\begin{criterion}[Promotion en univers descendant physique]
\label{crit:descendant-promotion}
L'image complète acquiert le statut de réalisation physique lorsque sont
construits simultanément :
\begin{enumerate}[label=\roman*)]
 \item une solution des équations Einstein--Cartan avec
       \cref{eq:descendant-bounce-conditions} ;
 \item des données d'Israel--Cartan régulières en \(\Sigma_H\) et \(\Sigma_b\) ;
 \item une région expansive future avec une fonction temporelle et des surfaces de
       Cauchy dans chaque branche accessible ;
 \item la finitude des invariants et le prolongement des géodésiques à travers le
       noyau ;
 \item l'absence de courbes causales fermées et de couches non enregistrées ;
 \item un canal isométrique qui conserve les charges déclarées et l'
       hamiltonien de provenance \cref{eq:descendant-HHM-intertwining} ;
 \item la naturalité nonadique du courant torsionnel et des données de
       frontière ;
 \item un entrelaceur entre le ruban généalogique et le fibré physique des
       repères ou des spineurs.
\end{enumerate}
\end{criterion}

Ces conditions transforment le récit
\cref{eq:descendant-five-stage-sequence} en un problème d'existence
géométrique, opératoriel et causal avec des contrôles séparés.

\section{Conséquences locales du raccordement torsionnel : conservation, régularité et prolongement causal}

\begin{proposition}[Contenu d'une réalisation complète]
\label{prop:descendant-conditioned-consequences}
Supposons satisfait le
\cref{crit:descendant-promotion}. Alors :
\begin{enumerate}[label=\alph*)]
 \item le signe causal parcourt la séquence
       \cref{eq:descendant-five-stage-sequence} ;
 \item la mémoire torsionnelle conserve la covariance \(V^{-2}\) des deux
       côtés du rebond ;
 \item l'état global évolue par un canal qui conserve la trace ;
 \item chaque observateur accède à un état relatif de sa branche ;
 \item les charges incluses dans
       \cref{eq:descendant-charge-energy-intertwining} sont conservées ;
 \item l'orientation de Möbius est transportée comme donnée d'holonomie.
\end{enumerate}
\end{proposition}

\begin{proof}
Chaque affirmation est l'application d'une des conditions du critère :
le lecteur causal prouve a) ; le transducteur torsionnel prouve b) ; l'isométrie
prouve c) et d) ; les entrelacements prouvent e) ; et le morphisme de fibrés prouve
f).
\end{proof}

La proposition organise les conséquences d'une structure construite. L'
existence de la continuation minimale est résolue par le
\cref{thm:ivv-descendant-existence} ; la ramification classique avec changement
topologique est classifiée par l'obstruction du
\cref{thm:ivv-cobordism-obstruction}.

\section{Dynamique torsionnelle du rebond : identités exactes, interfaces cosmologiques et conditions de clôture}

\begin{center}
\small
\begin{tabularx}{\textwidth}{>{\bfseries\raggedright\arraybackslash}p{31mm}YY}
\toprule
Objet & Statut & Obligation ou contrôle\\
\midrule
Loi \(M_n=V_n^{-2}\) & Résultat interne exact & Certificat de Perron et raffinement\\
Lecteur \(-\operatorname{sgn}\vartheta\) & Formalisation exacte & Observable d'expansion typé\\
Rebond homogène & Consecuencia condicionada & \(s_0>0\), \(w<1\), zéro transverse\\
Rebond Kantowski--Sachs & Interface physique & Dominance de la torsion, cisaillement et courbure contrôlés\\
Raccordement Israel--Cartan & Interface physique & Données de couche et courant de spin déclarés\\
Ruban de Möbius & Résultat topologique exact & Entrelaceur vers le fibré physique\\
Estado relativo & Résultat opératoriel exact & Instrument et fibre observée\\
Canal parent--descendant & Formalisation conditionnée & Isométrie et conservation des charges\\
Continuation descendante & Théorème d'existence & Solution régulière et causalement complète\\
Cobordismo ramificado & Obstruction classique démontrée & Incompatibilité avec l'hyperbolicité globale et le changement topologique\\
\bottomrule
\end{tabularx}
\end{center}

\section{Obstructions à la continuation cosmologique : perte d'orientation, de régularité du raccordement ou de mémoire de retour}

Les contrôles négatifs distinguent chaque couche :
\begin{enumerate}
 \item une solution avec \(\Theta=0\) et \(\dot\Theta<0\) représente un
       maximum de volume et ne relève pas du rebond ;
 \item le cas \(8\pi Gs_0\leq\Sigma_0^2\) élimine la garantie de
       dominance torsionnelle anisotrope ;
 \item un \(q_b\notin\Z\) conserve le rebond continu et réfute la version
       synchronisée avec un tour complet ;
 \item un saut de \(h_{ij}\), de la quantité d'Israel ou du courant de
       Cartan requiert une couche superficielle explicite ;
 \item un ruban avec \(w_1=0\) possède une orientation globale et réfute la
       réalisation de Möbius ;
 \item une application \(V_{\rm pc}\) avec \(V_{\rm pc}^*V_{\rm pc}\ne I\) perd de la
       probabilité globale ;
 \item deux copies identiques d'un état arbitraire violent le contrôle de non-
       clonage et réfutent la lecture de distribution corrélée ;
 \item une région descendante avec des courbes causales fermées, des invariants
       divergents ou des géodésiques inextensibles ne satisfait pas la promotion causale ;
 \item un changement topologique qui conserve simultanément toutes les
       hypothèses classiques faisant l'objet d'une obstruction est dépourvu du domaine géométrique
       requis.
\end{enumerate}

\begin{proofstatusbox}
La grammaire \(-1,0,+1\), le lecteur de signe, l'involution
\(C_M^2=I\), la classification du ruban, la loi \(V^{-2}\), la hiérarchie
des retours et la conservation d'un canal isométrique sont des résultats
exacts dans leurs domaines. Le rebond Einstein--Cartan est une conséquence des
équations et des inégalités déclarées. L'extension à
Kantowski--Sachs, les raccordements, la continuité de la connexion de Barbero et
l'ouverture causale sont des interfaces physiques. La continuation
parentale--descendante minimale est réalisée par le
\cref{thm:ivv-descendant-existence} ; le changement topologique ramifié sous
les hypothèses classiques est exclu par le
\cref{thm:ivv-cobordism-obstruction}.
\end{proofstatusbox}

\begin{falsifierbox}
La couture torsionnelle est réfutée si le courant construit à partir de HMT
produit une puissance différente de \(V^{-2}\), une amplitude de signe
incompatible ou un défaut de raffinement. Le rebond est réfuté par l'
absence d'une racine transverse régulière. Le raccordement est réfuté par un
résidu d'Israel ou de Cartan sans source superficielle. La conservation globale
est réfutée par une perte de trace, de norme ou d'une charge entrelacée.
La promotion en univers descendant est réfutée par une singularité,
une violation causale ou l'impossibilité de construire le fibré d'orientation et le
canal global exigés. Chaque réfutation affecte son interface et conserve les
résultats internes précédents.
\end{falsifierbox}
